\customchapter{ABSTRACT}
% Traduccion fiel de resumen.tex, correccion conjunta r02. Trazabilidad: redaccion/rondas/abstract-r02-respuesta.md.

\begin{center}
\large \vspace{-1.5cm} \textbf{Conditioned synthesis of transiliac-transsacral screws and their local metal artifacts in pelvic CT using image-domain diffusion with multi-window encoding}
\end{center}

Metallic implants produce streaks and dark regions in pelvic computed tomography that can obscure bone boundaries. Because scans with implants and bone annotations are scarce, generating images with implants from metal-free scans is proposed. This work proposed a pipeline that places a transiliac-transsacral screw and synthesizes its local metal artifact in implant-free scans from the CTPelvic1K dataset; its two components were evaluated separately.

First, it was tested whether a network that compresses and reconstructs images preserved the bone values needed for synthesis in a compressed space. Next, a placement sampler proposed three-dimensional screw positions, and a measure of surgical admissibility was defined from osseous corridor viability and cortical breach grade, which records how far the screw protrudes beyond the bone. Image synthesis used diffusion directly on image values with multi-window encoding; only the implant and a band of about 12~mm around it were regenerated. The intensity of the generated streaks, measured by \emph{streak amplitude}, was compared with a physics-based simulation protocol and with artifact-free screw insertion as a control.

The compression test failed: none of the six evaluated configurations met the preset limit of 25 Hounsfield units of bone error; the lowest mean error was 61.72 Hounsfield units in 34 patients. Consequently, synthesis was implemented directly on image values. In 72 patients, 3\,600 positions were generated; the distribution of their breach grades lay at distances of 0.206 and 0.230 grades from the two reference clinical series, separated by 0.395 grades. This closeness is descriptive because no criterion for satisfactory placement had been specified. In 13 patients, the mean intensity of the synthesized streaks was 301.0 Hounsfield units lower than that of the physics-based protocol with titanium. The 90\,\% confidence interval for this difference crossed the equivalence bound: equivalence of the two intensities could not be established.

The pipeline was implemented, but equivalence of its streak intensity with that of the physics-based protocol was not established. The author judged its visual realism insufficient; moreover, the synthesized artifact did not reproduce the decrease in streak intensity with distance from the metal. Whether these images improve the training of bone segmentation networks was not evaluated.

\noindent \textbf{Keywords:}\\
\noindent Medical image synthesis; Metal artifact; Pelvic computed tomography; Diffusion models; Iliosacral fixation
